\medskip

On considère l'équation différentielle 
\[\left(E_0\right) :\quad  y' = y\]
où $y$ est une fonction dérivable de la variable réelle $x$.

\medskip

\begin{enumerate}
\item Démontrer que l'unique fonction constante solution de l'équation différentielle 
$\left(E_0\right)$ est la fonction nulle.
\item Déterminer toutes les solutions de l'équation différentielle $\left(E_0\right)$.

On considère l'équation différentielle 
\[(E) :\quad  y' = y - \cos (x) - 3\sin (x)\]
où $y$ est une fonction dérivable de la variable réelle $x$.

\item La fonction $h$ est définie sur $\R$ par $h(x) = 2 \cos (x) + \sin (x)$.

On admet qu'elle est dérivable sur $\R$.

Démontrer que la fonction $h$ est solution de l'équation différentielle $(E)$.
\item On considère une fonction $f$ définie et dérivable sur $\R$.

Démontrer que : \og $f$ est solution de $(E)$ \fg{} est équivalent à \og $f - h$ est solution de $\left(E_0\right)$ \fg.
\item En déduire toutes les solutions de l'équation différentielle $(E)$.
\item Déterminer l'unique solution $g$ de l'équation différentielle $(E)$ telle que $g(0) = 0$.
\item Calculer:
\[\displaystyle\int_0^{\frac{\pi}{2}} \left[-2\e^x + \sin (x) + 2\cos (x)\right]\:\text{d}x.\]
\end{enumerate}

\medskip
